% ================================================================
% Drop-in section for notes.tex
% Requires \usepackage{amsmath,amssymb} and \usepackage{natbib}.
% Bibliography entries: advanced_stochastic_processes_readings_seven.bib
% ================================================================
\section{Stochastic Markov processes at work: seven selected papers}
\label{sec:seven-applications}

The following seven studies illustrate how stochastic descriptions of
individual events can help explain ecological and epidemiological patterns.
Together, they link continuous-time Markov processes to species
abundance distributions, mainland--island metapopulation dynamics,
nonlinear resource harvesting, biodiversity, epidemic oscillations,
spatial dispersal networks, and hybrid stochastic--deterministic systems. For each study, we identify the
scientific question, summarize its main ideas, and indicate the mathematical
concepts that students should understand before reading the paper.

\subsection{Kendall (1948): Population processes and Fisher's logseries}
\label{subsec:kendall}
\noindent\textbf{Reference.} \citet{Kendall1948Logseries}.

\medskip\noindent\textbf{Scientific question.}
Can elementary stochastic processes of reproduction, mortality, and
immigration provide a mechanistic explanation for Fisher's logarithmic series
species-abundance distribution?

\medskip\noindent\textbf{Main ideas.}
Kendall studies population growth models expressed through probabilities of
individual events occurring in continuous time. A useful starting point is the
linear immigration--birth--death process, with transition rates
\begin{equation}
  g_n=\nu+bn, \qquad r_n=dn,
\end{equation}
where $\nu$, $b$, and $d$ are immigration, per-capita birth, and per-capita
death rates, respectively. When $d>b$, this process admits a stationary
negative-binomial distribution. This should not be confused with the
logarithmic-series distribution of species abundances. The latter arises when
one considers the abundance distribution of immigrant-founded species or
lineages under suitable assumptions. Its normalized form is
\begin{equation}
  P(n)=-\frac{1}{\ln(1-x)}\,\frac{x^n}{n},
  \qquad n=1,2,\ldots,\quad 0<x<1.
\end{equation}
The paper establishes an early, influential connection between explicitly
stochastic population mechanisms and a recurring ecological pattern.

\medskip\noindent\textbf{Mathematical concepts to review.}
\begin{itemize}
  \item Linear birth--death--immigration processes and transition rates.
  \item Master equations for one-step Markov processes.
  \item Probability generating functions and stationary distributions.
  \item Negative-binomial and logarithmic-series distributions, and the
        distinction between population abundance and species abundance.
\end{itemize}

\noindent\textbf{Suggested task.} Derive the stationary distribution of the
linear immigration--birth--death process and explain what additional
interpretation or assumptions are needed to recover a logseries for species
abundances.

\subsection{Alonso and McKane (2002): Immigration and extinction in mainland--island metapopulations}
\label{subsec:alonso-mckane-metapopulations}
\noindent\textbf{Reference.} \citet{AlonsoMcKane2002Extinction}.

\medskip\noindent\textbf{Scientific question.}
How do immigration from a mainland, local colonization, and stochastic patch
extinction jointly determine the persistence and eventual extinction of a
metapopulation distributed over a finite number of habitat patches?

\medskip\noindent\textbf{Main ideas.}
The paper generalizes Levins' classical metapopulation model by introducing
immigration from an external mainland and explicitly treating the number of
occupied patches as a continuous-time stochastic variable. For $N$ patches,
let $n\in\{0,\ldots,N\}$ denote the number of occupied patches. Transitions
$n\to n+1$ describe immigration-driven or internal colonization of empty
patches; transitions $n\to n-1$ describe local extinctions. The corresponding
one-step master equation determines the time-dependent probability of
observing $n$ occupied patches. Its stationary solution characterizes patch
occupancy under persistent mainland immigration. The authors also analyze
fluctuations, the probability that all patches are empty, and mean times to
metapopulation extinction, using approximations and numerical simulations.
They then test how far the predictions carry over to lattice and spatially
realistic patch networks. An important subtlety is that an empty
metapopulation is not absorbing when mainland immigration remains possible;
the paper's extinction-time results must therefore be interpreted with
attention to the boundary conditions and immigration assumptions.

\medskip\noindent\textbf{Mathematical concepts to review.}
\begin{itemize}
  \item Finite-state birth--death processes for patch occupancy.
  \item Colonization and extinction rates, including external immigration.
  \item Master equations, boundary conditions, and stationary distributions.
  \item Transient extinction probabilities and mean first-passage times.
  \item Gaussian or linear-noise approximations and finite-size effects.
  \item Connections between well-mixed patch occupancy and spatial networks.
\end{itemize}

\noindent\textbf{Suggested task.} Write the transition rates for the
number of occupied patches and derive the corresponding master equation.
Explain the influence of mainland immigration on occupancy and extinction,
and distinguish a stationary probability of an empty metapopulation from an
absorbing extinction state.


\subsection{Cuenda, Llorente and Capit\'an (2020): Collapse and recovery under nonlinear harvesting}
\label{subsec:cuenda-llorente-capitan}
\noindent\textbf{Reference.} \citet{CuendaLlorenteCapitan2020Harvesting}.

\medskip\noindent\textbf{Scientific question.}
How does demographic stochasticity affect the timescales of collapse and
recovery of a renewable resource, such as a fish population, under nonlinear
harvesting and immigration? Can stochastic transition times differ
systematically from predictions of a deterministic population model?

\medskip\noindent\textbf{Main ideas.}
The authors study a resource subject to density-dependent population growth,
immigration, and nonlinear harvesting. In the deterministic description,
saturating extraction can generate alternative population regimes and
bifurcation thresholds associated with collapse and recovery. A related
continuous-time birth--death process keeps the population discrete and its
demographic events stochastic. Rather than asking only for the stationary
abundance distribution, the paper focuses on \emph{transition times}: how
long it takes the resource population to cross a threshold marking collapse
or recovery. Analytical and numerical comparisons show that stochastic
collapse and recovery times are bounded above by the corresponding
deterministic estimates under the conditions examined, across the range of
immigration rates considered. This illustrates that demographic noise can
accelerate transitions in both directions, rather than simply increasing
extinction risk.

For a generic one-step model with upward rates $g_n$ and downward rates
$r_n$, the mean first-passage time $T_n$ to a specified absorbing threshold
satisfies the backward equation
\begin{equation}
  -1=g_n(T_{n+1}-T_n)+r_n(T_{n-1}-T_n),
\end{equation}
with boundary conditions appropriate to the chosen threshold and population
state space. This is a general methodological illustration, not a
replacement for the paper's specific model and boundary conditions.

\medskip\noindent\textbf{Mathematical concepts to review.}
\begin{itemize}
  \item One-dimensional continuous-time birth--death processes with nonlinear rates.
  \item Logistic population growth, saturating harvesting, and deterministic bifurcations.
  \item Master equations and stochastic trajectories.
  \item First-passage times and backward equations.
  \item Boundary conditions and the distinction between extinction, collapse,
        and recovery thresholds.
  \item Comparison of stochastic and deterministic transition times.
\end{itemize}

\noindent\textbf{Suggested task.} Construct an event table for a population
with density-dependent growth, immigration, and saturating harvesting.
Derive its master equation and the backward recurrence for the mean time to
reach a collapse threshold. Discuss how the deterministic model and the
stochastic model characterize collapse and recovery differently.

\subsection{Etienne and Alonso (2007): Neutral community theory}
\label{subsec:etienne-alonso}
\noindent\textbf{Reference.} \citet{EtienneAlonso2007Neutral}.

\medskip\noindent\textbf{Scientific question.}
How can the stochastic balance of local immigration and extinction, together
with regional speciation and extinction, explain the coexistence of many
species without invoking species-specific niche differences?

\medskip\noindent\textbf{Main ideas.}
Neutral community theory describes species as demographically equivalent at
the level of the processes being modeled, while allowing stochastic birth,
death, dispersal, and speciation events to generate unequal abundances.
Dispersal limitation connects local community dynamics to a larger regional
species pool. The resulting theory provides predictions for diversity and
species-abundance patterns and clarifies the distinct roles of ecological
drift, immigration, and speciation. This study develops and synthesizes the
mathematical framework used to interpret neutral community models and their
sampling distributions.

\medskip\noindent\textbf{Mathematical concepts to review.}
\begin{itemize}
  \item Multispecies stochastic population processes and ecological drift.
  \item Birth--death--immigration dynamics in local communities.
  \item Neutrality, dispersal limitation, and local--regional coupling.
  \item Stationary distributions and sampling distributions.
  \item Probability generating functions and combinatorial probability,
        where relevant to the chosen derivation.
\end{itemize}

\noindent\textbf{Suggested task.} Select one neutral-community prediction
and explain which demographic events and dispersal assumptions lead to it.
Discuss how a local community differs from the regional species pool.

\subsection{Alonso, McKane and Pascual (2007): Stochastic amplification in epidemics}
\label{subsec:epidemics}
\noindent\textbf{Reference.} \citet{AlonsoMcKanePascual2007Epidemics}.

\medskip\noindent\textbf{Scientific question.}
Can demographic stochasticity sustain pronounced epidemic oscillations even
when the corresponding deterministic epidemic model has a stable equilibrium?

\medskip\noindent\textbf{Main ideas.}
The model describes discrete infection and demographic events, allowing for
infections imported from outside the focal population. Around a stable
endemic state, the intrinsic fluctuations generated by these events can be
amplified by the dynamical response of the system. Hence, stochastic
realizations may display relatively coherent epidemic oscillations although
the deterministic solution relaxes to equilibrium. A system-size expansion
provides an analytical approximation to fluctuations and their power spectrum.
This establishes a connection between demographic noise, dynamical
stability, and recurrent epidemics.

\medskip\noindent\textbf{Mathematical concepts to review.}
\begin{itemize}
  \item Multivariate continuous-time Markov jump processes.
  \item Master equations for epidemic compartment models.
  \item Deterministic limits and linear stability of equilibria.
  \item van Kampen's system-size expansion and linear noise approximation.
  \item Fourier transforms, covariance functions, and power spectra.
\end{itemize}

\noindent\textbf{Suggested task.} Explain, using the deterministic Jacobian
and the stochastic fluctuation equations, how stable deterministic dynamics
can amplify demographic noise into recurrent epidemics.

\subsection{Nicoletti et al. (2023): Effective metapopulation dynamics from dispersal}
\label{subsec:metapopulation}
\noindent\textbf{Reference.} \citet{NicolettiEtAl2023Metapopulation}.

\medskip\noindent\textbf{Scientific question.}
Can effective colonization rules and metapopulation persistence criteria be
derived from a mechanistic model that explicitly describes individual
dispersal across a spatial network?

\medskip\noindent\textbf{Main ideas.}
The authors distinguish settled individuals, which occupy habitat patches,
from dispersing individuals (explorers), which move through a network before
establishing in suitable patches. When dispersal dynamics are fast relative
to slower local changes, the explorer variables can be eliminated in a
reduced description. The resulting effective dispersal kernel reflects the
structure of the underlying network and yields generalized metapopulation
dynamics. In this way, effective colonization processes and persistence
thresholds arise from explicit microscopic mechanisms rather than being
specified solely at the macroscopic level.

\medskip\noindent\textbf{Mathematical concepts to review.}
\begin{itemize}
  \item Spatial Markov processes and network-based dispersal.
  \item State vectors, transition rates, and occupancy dynamics.
  \item Fast--slow dynamics and separation of timescales.
  \item Quasi-steady-state elimination and effective model reduction.
  \item Dispersal kernels, matrices, eigenvalues, and persistence thresholds.
\end{itemize}

\noindent\textbf{Suggested task.} Reconstruct the roles of settled
individuals and explorers, outline the fast-variable elimination, and
explain how dispersal-network topology becomes encoded in the reduced model.

\subsection{Realpe-Gomez, Galla and McKane (2012): Demographic noise in piecewise-deterministic systems}
\label{subsec:pdmp}
\noindent\textbf{Reference.} \citet{RealpeGomezGallaMcKane2012PDMP}.

\medskip\noindent\textbf{Scientific question.}
How can a large population of discrete individuals interact with continuous
environmental variables, and how does intrinsic demographic noise change the
qualitative dynamics, including the emergence of oscillations?

\medskip\noindent\textbf{Main ideas.}
The authors study hybrid dynamical systems in which the state includes a
population of discrete individuals and environmental variables evolving
continuously. Between discrete demographic events, environmental variables
follow deterministic evolution equations; demographic events occur at
state-dependent stochastic rates and alter the subsequent dynamics. This
provides a mechanistic route to piecewise-deterministic Markov processes.
The framework is motivated in part by semi-arid vegetation systems, where
plants are discrete and water availability is continuous. Through a
system-size expansion and numerical analysis, the paper quantifies
fluctuations around macroscopic dynamics and derives their power spectra.
These calculations show that demographic noise can generate pronounced,
characteristic oscillations absent from the corresponding deterministic
trajectory. The central contribution is the theoretical treatment of
stochastic hybrid dynamics, rather than the detailed claims about vegetation
resilience developed in the subsequent 2013 paper.

\medskip\noindent\textbf{Mathematical concepts to review.}
\begin{itemize}
  \item Continuous-time Markov jump processes and state-dependent hazards.
  \item Piecewise-deterministic Markov processes and hybrid state spaces.
  \item Deterministic flow between stochastic jump events.
  \item Master or forward equations combining jumps and continuous transport.
  \item System-size expansions and Gaussian fluctuation approximations.
  \item Power spectral density and noise-induced quasi-cycles.
\end{itemize}

\noindent\textbf{Suggested task.} Define a minimal hybrid state consisting
of one discrete population variable and one continuous environmental
variable. Describe the event rates, the deterministic flow between jumps,
and the key steps through which a system-size expansion predicts the
spectrum of noise-induced fluctuations.

\subsection{A common framework for discussing the seven papers}
\label{subsec:common-framework}

Each student's presentation should address the same five questions:
\begin{enumerate}
  \item What ecological or epidemiological question motivates the model?
  \item What are the state variables, elementary events, and transition rates?
  \item Which mathematical tools are used to analyze or approximate the
        stochastic process?
  \item What is the principal scientific result?
  \item What would be missed, altered, or obscured by considering only a
        deterministic description?
\end{enumerate}

Where feasible, students should present a table of elementary events and
rates, a schematic of the stochastic mechanisms, and one central figure or
analytical result from their assigned paper. Across the seven contributions,
the unifying methodological idea is to begin with individual-level events
and derive probabilistic and, when justified, macroscopic descriptions.
